\section*{Completion of the two-bit half-tax continuum}

All logarithms in this proof are natural.  For $0<p\le1/2$ and
$0\le q\le1$, put
\[
 s=4q(1-q),\quad c=4q^2,\quad b=1-sp,\quad
 a={cp(1-p)\over b},
\]
and
\[
 d(p,q)=b h(a)+sp h(p)-h(sp)+p h(s)
          +{p\over2}\log{1-p\over p},                 \tag{1}
\]
where $h(x)=-x\log x-(1-x)\log(1-x)$.  Thus
$d=(\ln2)(M_A+\Delta_A/2-J_A)$ for the frozen two-bit Boolean AND
kernel.

\paragraph{Theorem.}
For every $0<p\le1/2$ and every $0\le q\le1$,
\[
 \boxed{d(p,q)\ge0.}                                  \tag{2}
\]
Equality on the two endpoint collars occurs only at
$(p,q)=(1/2,0)$ and $(1/2,1)$.  Together with the previously certified
interior and rare-bias regions, this proves the half-tax inequality on
the complete two-parameter domain.

\paragraph{Balanced lower edge.}
At $p=1/2$, direct entropy algebra gives
\[
 \begin{aligned}
 d(1/2,q)={}&-2q^2\log q-2(1-q)^2\log(1-q)\\
 &+2\{q^2+(1-q)^2\}\log\{q^2+(1-q)^2\}
 -(1-2q)^2\log|1-2q|.                         \tag{3}
 \end{aligned}
\]
For $0\le q\le1/100$, define
\[
 g(x)=-(1-x)^2\log(1-x),\qquad
 y=2q(1-q).
\]
The convergent power series give
\[
 g(x)=x-\frac32x^2+
 \sum_{n\ge3}{2x^n\over n(n-1)(n-2)}
 \ge x-\frac32x^2,                                   \tag{4}
\]
and
\[
 (1-y)\log(1-y)=-y+
 \sum_{n\ge2}{y^n\over n(n-1)}\ge-y+\frac12y^2.     \tag{5}
\]
Rewriting (3) as
\[
 d(1/2,q)=-2q^2\log q+2g(q)+g(2q)
                 +2(1-y)\log(1-y)
\]
and applying (4)--(5) yields
\[
 d(1/2,q)\ge q^2\left(2\log{1\over q}-1-8q+4q^2\right)>0       \tag{6}
\]
for $0<q\le1/100$, while $d(1/2,0)=0$.

\paragraph{Monotonicity away from the balanced edge.}
Differentiating (1) is legitimate in the open collar.  The potentially
singular products are evaluated in the continuous forms
\[
 b h'(a)a'={cR\over b}
  \{\log m-\log c-\log p-\log(1-p)\},                 \tag{7}
\]
where
\[
 R=(1-2p)b+sp(1-p),\qquad m=b-cp(1-p),
\]
and
\[
 -s h'(sp)=-s\log(1-sp)+s\log s+s\log p.              \tag{8}
\]
The conventions $x\log x=0$ at $x=0$ make (7)--(8) continuous on the
closed rectangle.

The accompanying directed-rounding certificate encloses this derivative
on
\[
 {49\over100}\le p\le{1\over2},\qquad
 0\le q\le{1\over100}.
\]
It proves the uniform bound
\[
 \boxed{\partial_p d(p,q)\le-0.9.}                    \tag{9}
\]
The largest certified upper endpoint is
\[
 -0.9002260317012236411683888873233463387246104808513.
\]
The proof uses 90-decimal directed arithmetic, visits 45 rectangles,
accepts 23 leaves, reaches depth 6, and leaves no unresolved rectangle.
No floating-point sign decision is made.  Integrating (9) and using (6)
gives
\[
 d(p,q)\ge d(1/2,q)+0.9(1/2-p)\ge0                  \tag{10}
\]
on the complete lower endpoint collar.

\paragraph{Upper endpoint collar without a second certificate.}
Fix $0\le q\le1/100$ and compare $q$ with $1-q$.  The values of $s$ and
$b$ are unchanged; only the first entropy term in (1) changes.  Its two
arguments are
\[
 a_-={4q^2p(1-p)\over b},\qquad
 a_+={4(1-q)^2p(1-p)\over b}.
\]
They satisfy $a_-\le a_+$ and
\[
 1-a_--a_+
 ={(1-2p)(1-2p+sp)\over b}\ge0.                      \tag{11}
\]
Thus $a_-\le a_+\le1-a_-$.  Since binary entropy is symmetric and
increases on $[0,1/2]$, $h(a_+)\ge h(a_-)$.  All other terms agree, so
\[
 d(p,1-q)\ge d(p,q).                                  \tag{12}
\]
Equation (10) therefore proves the entire upper collar
$99/100\le q\le1$ as well.

\paragraph{Complete coverage and exact GU implication.}
The source-bound pieces are now:
\[
\begin{array}{c|c}
0<p\le1/64 & \text{normalized rare-tail theorem},\\
1/64\le p\le49/100 & \text{repaired directed certificate, all }q,\\
49/100\le p\le1/2,\ 1/100\le q\le99/100
 & \text{balanced-transition and central-collar certificates},\\
49/100\le p\le1/2,\ q\in[0,1/100]\cup[99/100,1]
 & \text{(6)--(12).}
\end{array}
\]
Hence (2) has no omitted bias or endpoint parameter.

This is still a two-bit local lemma, not universal GU.  If its mismatch
terms $M_A$ vanished globally and every $\Delta_A/2$ were charged once,
one would obtain $G\le I/2$.  The growing-rank OR-defect family has
$G/I\to1$ (indeed $G>I$ eventually), so those global premises cannot both
hold.  The completion therefore localizes the remaining debt: a universal
proof must retain a nonzero higher-arity, hierarchy, or terminal
$M$-reserve.  No claim that the half-tax lemma alone proves GU or FEI is
made.
